\chapter{Related Work}
\label{chap:relatedwork}

This chapter covers four bodies of work. Section~\ref{sec:rw_kgs} introduces the
knowledge graphs this thesis measures. Section~\ref{sec:rw_quality_metrics}
reviews metrics that describe one snapshot, and
Section~\ref{sec:rw_ranking} the ranking measures that
Chapter~\ref{chap:approach} builds on. Section~\ref{sec:rw_evolution} covers work
on how knowledge graphs change, and Section~\ref{sec:rw_comparison} the
comparison of versions and of different graphs.
Section~\ref{sec:rw_gaps} states what none of it provides.

\section{Knowledge Graphs and Their Releases}
\label{sec:rw_kgs}

A knowledge graph stores facts as RDF triples~\cite{cyganiak2014rdf} and exposes
them to queries through SPARQL~\cite{harris2013sparql}. Hogan et
al.~\cite{hogan2021knowledge} survey the field and the design decisions that
separate one graph from another.

The two graphs measured here have long release histories. YAGO began by combining
Wikipedia categories with WordNet~\cite{suchanek2007yago,miller1995wordnet}, gained spatial and
temporal facts in YAGO2~\cite{hoffart2013yago2}, and became multilingual in
YAGO3~\cite{mahdisoltani2015yago3}. YAGO~4 broke with all of
that~\cite{tanon2020yago}: it dropped the WordNet and Wikipedia-category
taxonomy and adopted schema.org~\cite{guha2016schema} classes with Wikidata
identifiers. YAGO~4.5~\cite{suchanek2024yago45} keeps the schema.org top level
and adds a large part of the Wikidata taxonomy beneath it, subject to logical
constraints that keep classes and instances apart, which is why its class count
grows from about ten thousand to over a hundred thousand. That break
matters to this thesis in a concrete way, because it means YAGO~3 and YAGO~4
share almost no class names, and Chapter~\ref{chap:experiments} reports what
happens when a per-class metric meets that discontinuity.

DBpedia extracts structured data from Wikipedia
infoboxes~\cite{auer2007dbpedia,bizer2009dbpedia}, and the release used
here is described by Lehmann et al.~\cite{lehmann2015dbpedia}. Since 2020
DBpedia has been published through a redesigned, largely automated release
cycle with monthly snapshots~\cite{hofer2020dbpedia}, which is the source of the
2022 and 2025 releases measured here. Wikidata is the
third large open graph~\cite{vrandecic2014wikidata,erxleben2014introducing};
this thesis does not measure it, and Section~\ref{sec:future_work} says why that
is the obvious next step.

All queries here run through QLever~\cite{bast2017qlever}, a SPARQL engine that
answers from a precomputed index and can hold a billion-triple graph on a single
machine. That capability is what makes the endpoint-based approach of
Chapter~\ref{chap:approach} possible at all.

\section{Quality and Importance Metrics}
\label{sec:rw_quality_metrics}

Zaveri et al.~\cite{zaveri2016quality} survey linked data quality assessment and
organise it into dimensions such as completeness, consistency and conciseness.
Färber et al.~\cite{farber2018linked} apply a comparable framework to DBpedia,
Freebase, OpenCyc, Wikidata and YAGO side by side, which is the closest prior
work to the cross-graph comparison in Section~\ref{sec:exp5}. Their comparison is
a one-off study of a fixed set of releases rather than a repeatable measurement,
and it does not track a graph across its own versions.

Seo et al.~\cite{seo2024structural} propose structural quality metrics computed
from the graph topology. Their measures describe one snapshot and say nothing
about a successor.

Dataset profiling tools share that limitation. LODStats~\cite{auer2012lodstats}
computes statistics such as class and property usage over a stream of RDF
statements, and Duan et al.~\cite{duan2011apples} define structuredness, a
measure of how uniformly the instances of a class use the properties of that
class, to show that benchmark data looks nothing like real RDF data. Both are
close in spirit to the class-level metrics of this thesis, and structuredness in
particular asks a question related to class entropy. Both describe one dataset
at a time.

The direct ancestor of this work is Knowgly, by García and
Bobed~\cite{garcia2025information}. Knowgly measures how much information a
predicate carries for a class, using property entropy and entropy-weighted type
importance, and it establishes the pattern this thesis reuses: SPARQL performs
the counting, the host language holds the aggregated counts and applies the
formula. Chapter~\ref{chap:approach} adopts that pattern and two of its metrics
as baselines. What Knowgly does not do is compare two graphs. Every Knowgly
number describes one endpoint, and nothing in it addresses what changes between
releases.

Shannon's definition of entropy~\cite{shannon1948mathematical} underlies metrics
2, 3, 4 and 7. Section~\ref{subsec:count_of_counts} shows how to compute it exactly
over a billion-triple graph without streaming the values.

\section{Entity and Predicate Ranking}
\label{sec:rw_ranking}

PageRank~\cite{page1999pagerank} ranks nodes by the stationary distribution of a
random walk. Applied to RDF without modification it treats every edge alike, so
a predicate that appears on every entity contributes as much as one that appears
rarely. On YAGO~4 this is a real problem: \texttt{schema:gender} carries a
normalised property entropy of 0.034, while \texttt{schema:geo} carries 1.000,
and unweighted PageRank cannot tell them apart.

Menendez et al.~\cite{menendez2018ranking} decouple ranking from raw
degree and rank entities by the information they carry, which motivates metric~5
in this thesis. Section~\ref{sec:entropy_weighted_pagerank} takes the next step
and pushes the information measure into the edge weights of the walk itself.

Weighting PageRank is not new in itself. Xing and Ghorbani~\cite{xing2004weighted}
distribute rank in proportion to the in-link and out-link popularity of the
target pages, and Langville and Meyer~\cite{langville2004deeper} survey the
variants and the conditions under which the power iteration converges. What
differs here is the source of the weight: it comes from the predicate that
labels the edge, measured by its entropy over the whole graph, rather than from
the link structure around the node.

\section{Knowledge Graph Evolution}
\label{sec:rw_evolution}

Polleres et al.~\cite{polleres2023knowledge} describe how knowledge evolves in
open knowledge graphs and argue that consumers need to understand change rather
than only consume the latest release. The ISWS report on evolution and
preservation~\cite{isws2019evolution} makes a similar case from the archival
side.

Käfer et al.~\cite{kafer2013observing} measured linked data dynamics empirically
by monitoring documents over time and reporting how often they change. Their unit
of observation is the document, not the class or the predicate, so the results
say how much a dataset churns without saying which part of the schema moved.

Ontology change has its own literature. Flouris et
al.~\cite{flouris2008ontology} classify the kinds of change an ontology can
undergo, and Noy and Musen~\cite{noy2004ontology} address versioning inside an
ontology management framework. Both work at the schema level and assume the
schema is small enough to reason about directly, which is true of a hand-built
ontology and false of DBpedia's 54{,}933 predicates.

\section{Comparing Versions and Comparing Graphs}
\label{sec:rw_comparison}

Papavasileiou et al.~\cite{papavasileiou2013high} detect high-level changes in
RDF(S) knowledge bases, lifting low-level triple additions and deletions into
meaningful operations. Roussakis et al.~\cite{roussakis2015flexible} generalise
this into a framework for defining and detecting changes over evolving RDF
datasets. Zeginis et al.~\cite{zeginis2011deltas} study how to compute the
delta between two RDF/S knowledge bases and how inferred triples change its
size. All three operate over datasets held locally, and all assume the two
versions can be loaded and traversed together.

The closest relative of metric~10 is the DBpedia Wayback
Machine~\cite{fernandez2015dbpedia}, which reconstructs historical DBpedia states
and diffs them using integer-encoded triples. Replacing IRIs and literals by
integer identifiers through a dictionary is standard practice in RDF storage,
from the RDF-3X engine~\cite{neumann2010rdf3x} to the HDT exchange
format~\cite{fernandez2013hdt}. This thesis takes the same encoding
idea and pairs it with a scoping scheme that removes the need for a global sort,
which is what allows the diff to run against a live endpoint rather than a
prepared archive.

The closest relative of metric~13 is Ringler and
Paulheim~\cite{ringler2017one}, who ask whether DBpedia, YAGO and Wikidata are
interchangeable and conclude that they are not. Their comparison is manual and
fixed in time. This thesis automates the class matching and reports the same
comparison as a repeatable metric, at the cost of matching classes by local name
rather than by instance-level alignment. Section~\ref{sec:exp5} is explicit about
what that costs.

\section{Summary of Gaps}
\label{sec:rw_gaps}

Three gaps follow from the above, and this thesis addresses each.

\textbf{Per-snapshot metrics are not tracked across versions.} Knowgly, Seo et
al.\ and the quality surveys all describe a single graph. A value becomes far
more informative when read next to the same value from the previous release, and
no existing tool computes a metric catalogue across a release series.

\textbf{Diffing work assumes local dumps.} Papavasileiou et al., Roussakis et
al., Zeginis et al.\ and the DBpedia Wayback Machine all operate on data the tool
controls. None
of them answers what an analyst can compute when the graph is only reachable
through a SPARQL endpoint, which is the normal situation for a billion-triple
graph on ordinary hardware.

\textbf{No combined metric and visual analytics pipeline exists.} Visual
analytics couples automatic analysis with interactive
visualisation~\cite{keim2008visual}, and it has reached large knowledge graphs:
the SHACL Dashboard of Mäkelburg et al.~\cite{makelburg2025shacl} makes
validation reports over large graphs explorable where reading them by hand is
infeasible. It analyses the quality of one graph against its shapes, not the
difference between releases. Comparison results are otherwise reported as
tables in papers. Nothing takes a metric catalogue across
several versions of several graphs and presents it so the movement is visible.
Chapter~\ref{chap:experiments} shows that the presentation is not cosmetic: the
same blank cell in a comparison can mean a measured zero, an unmeasured class or
a class that does not exist in that release, and a reader who cannot tell them
apart will draw the wrong conclusion.
